\documentclass[11pt]{article}
\usepackage[T1]{fontenc}
\usepackage{lmodern}
\usepackage{amsmath,amssymb,amsthm}
\usepackage[margin=1.05in]{geometry}
\usepackage[hidelinks]{hyperref}
\hypersetup{pdftitle={Continuity of states for bounded admissible inputs},
  pdfauthor={Kaiwen Jin},pdfsubject={AIM 411: bounded admissible inputs}}
\newtheorem{theorem}{Theorem}
\newtheorem{lemma}[theorem]{Lemma}
\theoremstyle{remark}
\newtheorem{remark}[theorem]{Remark}
\newcommand{\Lin}{\mathcal L}
\newcommand{\R}{\mathbb R}
\newcommand{\C}{\mathbb C}
\newcommand{\N}{\mathbb N}
\newcommand{\ess}{\operatorname*{ess\,sup}}
\title{Continuity of states for bounded admissible inputs}
\author{Kaiwen Jin\\ \small\texttt{yui.ui.siki@gmail.com}}
\date{October 15, 2026}
\begin{document}
\maketitle
\begin{abstract}
Let \(A\) generate a strongly continuous semigroup on an arbitrary Banach
space \(X\), and let \(B\) be a bounded operator from an arbitrary Banach
input space into the extrapolation space \(X_{-1}\).
We prove that \(L^\infty\)-admissibility for all strongly measurable inputs
implies that the finite-time control operator norms tend to zero with the
time horizon. Consequently every state driven by a bounded strongly
measurable input is continuous in the norm of \(X\), including at time zero.
The argument applies the scalar Grothendieck property to a separate
multiplier operator for each fixed input. A resolvent construction gives
this operator a separable range in \(X\); the Orlicz--Pettis theorem and a
disjoint-interval construction then give the uniform small-time estimate.
\end{abstract}

\section{The continuity question}

For an unbounded control operator, admissibility puts each convolution
state in the original state space, although the integral is initially
defined in an extrapolation space. It does not directly give continuity
of the state in the stronger norm. Weiss posed the bounded-input
continuity question in \cite[Problem~2.4]{Weiss}.
Preu{\ss}ler and Schwenninger obtained positive answers under several
geometric assumptions on the state and input spaces
\cite[Corollary~4.10 and the discussion following Corollary~4.11]{PS}.
We prove the unrestricted Banach-space assertion for the extrapolation
control-operator formulation.

Let \(X,U\) be Banach spaces over the same field, either \(\R\) or \(\C\).
Let \(A:D(A)\subset X\to X\) generate a \(C_0\)-semigroup \(S\).
Fix a resolvent point \(\lambda\) and form the completion \(X_{-1}\) of
\(X\) in the norm
\[
 \|x\|_{-1}=\|(\lambda-A)^{-1}x\|_X .
\]
Write \(J:X\to X_{-1}\) for the continuous injective canonical map and
\(S_{-1}\) for the compatible \(C_0\)-semigroup on \(X_{-1}\).
The resolvent extends to a bounded map
\(R_{-1}:X_{-1}\to X\). Its restriction to \(X\) is
\(R=(\lambda-A)^{-1}:X\to D(A)\), and
\[
 R_{-1}S_{-1}(t)=S(t)R_{-1}.
\]
These are the usual extrapolation constructions; see
\cite[Definition~3.1, Proposition~3.3 and Remark~3.4]{Weiss}.
We identify \(X\) with \(J(X)\) in convolution formulas.

Throughout, \(L^\infty(0,T;U)\) consists of equivalence classes of
\emph{strongly measurable}, essentially bounded functions, with the
essential supremum norm. All measures are completed Lebesgue measures,
and all integrals are Bochner integrals, initially in \(X_{-1}\).
For \(B\in\Lin(U,X_{-1})\), assume that for every \(T>0\) there is
\(K_T<\infty\) such that
\begin{equation}\label{eq:admissibility}
 \Phi_Tu:=\int_0^T S_{-1}(T-s)Bu(s)\,ds\in X,\qquad
 \|\Phi_Tu\|_X\le K_T\|u\|_\infty
 \quad(u\in L^\infty(0,T;U)).
\end{equation}
Time reversal gives the bounded operator
\begin{equation}\label{eq:reversed}
 F_hv=\int_0^h S_{-1}(r)Bv(r)\,dr,\qquad
 \kappa(h)=\|F_h\|_{\Lin(L^\infty(0,h;U),X)} .
\end{equation}
Thus \(F_h\) and \(\Phi_h\) have the same norm. Zero extension in the
reversed-time formula gives \(\kappa(h)\le\kappa(T)\) for \(0<h\le T\).

\begin{theorem}\label{thm:main}
Under \eqref{eq:admissibility}, with no further assumptions on \(X\) or
\(U\), one has
\begin{equation}\label{eq:zero}
 \lim_{h\downarrow0}\kappa(h)=0.
\end{equation}
For every \(T>0\) and \(u\in L^\infty(0,T;U)\), the function
\[
 x_u(0)=0,\qquad
 x_u(t)=\int_0^t S_{-1}(t-s)Bu(s)\,ds\quad(0<t\le T)
\]
belongs to \(C([0,T];X)\).
\end{theorem}

The norm \(\kappa(h)\) in \eqref{eq:zero} is taken over the entire strongly
measurable input unit ball. Its decay is stronger than continuity for one
fixed input and is called zero-class admissibility.
To obtain it, we first fix an input \(w\) and vary only a scalar multiplier
of \(w\). The resulting operator has separable range even when neither
\(X\) nor \(U\) is separable. This allows the scalar \(L^\infty\)
Grothendieck theorem to be used without assuming a Grothendieck property
for vector-valued \(L^\infty\).

\section{A vector measure for each fixed input}

We use two Banach-space results. Every bounded operator from scalar
\(L^\infty(0,T)\) into a separable Banach space is weakly compact:
scalar \(L^\infty\) is Grothendieck, and the asserted mapping property is
one of the characterizations of that property
\cite[Theorem~3.1.1 and Proposition~4.1.5]{GK}.
The Orlicz--Pettis theorem says that a series in a Banach space whose
every subseries converges weakly has every subseries convergent in norm
\cite[Chapter~IV, p.~24]{Diestel}.
Here weak convergence always uses the full continuous dual.

These statements can be used over either scalar field. In particular,
for a real Banach space \(E\) use its complexification with norm
\[
 \|x+iy\|_{E_\C}
   =\sup_{\theta\in\R}\|\cos\theta\,x-\sin\theta\,y\|_E .
\]
A bounded real operator \(Q:L^\infty(0,T;\R)\to E\) extends boundedly
to complex scalar \(L^\infty\) with values in \(E_\C\).
Separability is preserved, and restriction followed by the real-part
projection transfers weak compactness back to \(Q\).
The same embedding transfers the Orlicz--Pettis statement to real
spaces: each complex functional restricts to real and imaginary
continuous real functionals.

\begin{lemma}\label{lem:range}
Fix \(T>0\) and \(w\in L^\infty(0,T;U)\). The scalar multiplier operator
\[
 Q_w:L^\infty(0,T)\to X,\qquad Q_wf=F_T(fw),
\]
has its range in a closed separable subspace of \(X\) and is weakly compact.
\end{lemma}

\begin{proof}
Admissibility gives
\(\|Q_wf\|_X\le K_T\|w\|_\infty\|f\|_\infty\).
Strong measurability gives a separable closed subspace \(U_0\subset U\)
containing the values of a representative of \(w\) outside a null set.
Indeed, take the closed span of the finite ranges of simple functions
converging to \(w\) almost everywhere. Change \(w\) on that null set if
necessary. Put \(C=R_{-1}B\in\Lin(U,X)\) and
\[
 E=\overline{\operatorname{span}}^{\,X}
       \{S(t)Ca:t\ge0,\ a\in U_0\}.
\]
This subspace is separable: rational times, a countable dense subset of
\(U_0\), and rational scalar coefficients give a dense countable set.
The density follows from strong continuity and the local operator bounds
of a \(C_0\)-semigroup. The semigroup law also shows that \(E\) is
\(S\)-invariant.

For \(x=Q_wf\in X\), bounded resolvent commutation with the integral gives
\begin{equation}\label{eq:resolvent}
 Rx=\int_0^T S(r)Cw(r)f(r)\,dr\in E .
\end{equation}
For a simple input the kernel is a finite combination of continuous
semigroup orbits. Simple approximants to the input, together with local
semigroup bounds, give strong measurability of the integrand.
The bounded input also gives an integrable constant bound on its norm.
Since the integrand is \(E\)-valued, its integral is in the closed subspace
\(E\). The same argument applies to the subsequent kernels.
Let \(y=Rx\). Then \(y\in D(A)\cap E\). Invariance implies
\((S(h)y-y)/h\in E\) for \(h>0\). Taking its \(X\)-norm limit gives
\(Ay\in E\), hence \(x=\lambda y-Ay\in E\).
Thus \(Q_w\) is bounded with values in the separable Banach space \(E\).
The scalar Grothendieck result makes it weakly compact into \(E\), and
the inclusion \(E\to X\) preserves weak compactness.
\end{proof}

\begin{lemma}\label{lem:measure}
For fixed \(T,w\) as in Lemma~\ref{lem:range}, the set function
\[
 m_w(D)=Q_w\mathbf1_D
       =\int_D S_{-1}(r)Bw(r)\,dr\in X
 \qquad(D\subset(0,T)\text{ measurable})
\]
is countably additive in the norm of \(X\).
Consequently, for any disjoint measurable \(D_n\),
\begin{equation}\label{eq:measure-limits}
 \|m_w(D_n)\|_X\longrightarrow0,\qquad
 \|m_w((0,\delta))\|_X\longrightarrow0\quad(\delta\downarrow0).
\end{equation}
\end{lemma}

\begin{proof}
Linearity gives finite additivity. Put
\(g(r)=S_{-1}(r)Bw(r)\in L^1(0,T;X_{-1})\).
For disjoint \(D_n\) and any subset \(I\subset\N\), let
\[
 z_N=\sum_{\substack{n\in I\\n\le N}}m_w(D_n),\qquad
 D_I=\bigcup_{n\in I}D_n .
\]
Each \(z_N\) is \(Q_w\) of an indicator of norm at most one.
By Lemma~\ref{lem:range}, all these vectors lie in a fixed weakly
compact subset \(K\) of \(X\). Dominated convergence in \(X_{-1}\) gives
\begin{equation}\label{eq:injected}
 Jz_N\longrightarrow\int_{D_I}g(r)\,dr
             =Jm_w(D_I)\quad\text{in }X_{-1}.
\end{equation}
If a subnet of \(z_N\) converges weakly in \(X\), weak continuity of \(J\)
and \eqref{eq:injected} identify its limit with \(m_w(D_I)\), since
\(J\) is injective. Compactness now gives weak convergence of the whole
sequence to that vector. Otherwise a weak neighborhood of \(m_w(D_I)\)
would exclude infinitely many terms; those terms would have a cluster
point in the closed complement within \(K\), contradicting this
identification.

This holds for every \(I\), so every subseries of
\(\sum_n m_w(D_n)\) converges weakly in \(X\).
Orlicz--Pettis gives norm convergence. Applying \(J\) and
\eqref{eq:injected} identifies the norm sum with \(m_w(D_I)\).
This proves norm countable additivity, and the first limit in
\eqref{eq:measure-limits} follows.
For the second, take any decreasing \(\delta_k\downarrow0\).
The intervals \((\delta_{k+1},\delta_k]\) partition
\((0,\delta_1]\), so \(m_w((0,\delta_k])\) are tails of a norm-convergent
series and tend to zero. Null endpoints have measure zero because their
injected integrals vanish and \(J\) is injective.
If the unrestricted \(\delta\)-limit failed, a decreasing violating
sequence would contradict the conclusion just proved.
\end{proof}

\section{The small-time estimate and continuity}

\begin{proof}[Proof of Theorem~\ref{thm:main}]
Suppose that \(\kappa(h)\) fails to tend to zero.
By the definition of operator norm there is \(\varepsilon>0\) such that
arbitrarily small \(h>0\) admit \(v\in L^\infty(0,h;U)\) with
\(\|v\|_\infty\le1\) and \(\|F_hv\|_X\ge\varepsilon\).
No attainment of the operator norm is needed.
Choose such \(h_1<1,v_1\).
Lemma~\ref{lem:measure}, applied to this one input on \((0,h_1)\),
gives \(0<\delta_1<h_1\) such that
\[
 \left\|\int_0^{\delta_1}S_{-1}(r)Bv_1(r)\,dr\right\|_X
       <\frac{\varepsilon}{2}.
\]
Inductively choose \(h_{n+1}<\delta_n/2\) with a corresponding \(v_{n+1}\),
and then \(0<\delta_{n+1}<h_{n+1}\) with the same small-tail estimate.
The intervals \(I_n=(\delta_n,h_n)\) are disjoint, and
\begin{equation}\label{eq:lower}
 y_n=\int_{I_n}S_{-1}(r)Bv_n(r)\,dr\in X,\qquad
 \|y_n\|_X\ge\frac{\varepsilon}{2}.
\end{equation}
Define \(w:(0,1)\to U\) to equal \(v_n\) on \(I_n\) and zero elsewhere.
Each input has essentially separable range; the closed span of their
countable union is separable. The finite pasted functions converge
pointwise to \(w\), so \(w\) is strongly measurable, for example by
the Pettis measurability theorem \cite[Chapter~IV, p.~25]{Diestel}.
A countable union of null exceptional sets suffices to arrange
\(\|w\|_\infty\le1\).
The original admissibility hypothesis applies to this single input and
every input \(fw\) with scalar \(f\in L^\infty(0,1)\).
Lemma~\ref{lem:measure} at horizon \(1\) therefore gives
\(m_w(I_n)=y_n\to0\) in \(X\), contradicting \eqref{eq:lower}.
This proves \eqref{eq:zero}.

Fix now \(T>0\), \(u\in L^\infty(0,T;U)\), and
\(M=\|u\|_\infty\). Admissibility puts every \(x_u(t)\) in \(X\), and
\[
 \|x_u(t)\|_X\le\kappa(t)M\longrightarrow0\quad(t\downarrow0).
\]
For \(t_0\in(0,T]\), choose \(0<d<t_0\) and set \(a=t_0-d\).
Splitting the convolution at \(a\), using the semigroup law and the
compatibility with \(X\), gives for \(a<t\le T\)
\[
 x_u(t)=S(t-a)x_u(a)
            +F_{t-a}\bigl(r\mapsto u(t-r)\bigr).
\]
For \(|t-t_0|<d/2\) with \(t\in[0,T]\), monotonicity of \(\kappa\) gives
\begin{equation}\label{eq:continuity}
 \|x_u(t)-x_u(t_0)\|_X
 \le \|[S(t-a)-S(t_0-a)]x_u(a)\|_X+2\kappa(2d)M .
\end{equation}
Given a positive tolerance, first choose \(d\) to make the second term
small. With \(d,a\) and \(x_u(a)\) now fixed, strong continuity of the
single semigroup orbit makes the first term small as \(t\to t_0\).
This proves both sides of the limit at interior times and the relative
left limit at \(T\). Together with the limit at zero, it proves
\(x_u\in C([0,T];X)\).
\end{proof}

\begin{remark}
The separable subspace in Lemma~\ref{lem:range} may depend on \(w\);
no common separable subspace is required for all inputs.
Weak compactness is used only for \(Q_w\), with no such assertion for
\(F_T:L^\infty(0,T;U)\to X\).
The theorem concerns the Bochner convolution with the given \(B\).
Input maps not assumed to have this representation are outside its statement.
\end{remark}

\begin{thebibliography}{9}
\bibitem{Diestel}
J.~Diestel, \emph{Sequences and Series in Banach Spaces},
Graduate Texts in Mathematics, vol.~92, Springer-Verlag, New York, 1984.
\href{https://doi.org/10.1007/978-1-4612-5200-9}{doi:10.1007/978-1-4612-5200-9}.

\bibitem{GK}
M.~Gonz{\'a}lez and T.~Kania,
\emph{Grothendieck spaces: the landscape and perspectives},
arXiv:2102.03838v3, 2021.
\url{https://arxiv.org/abs/2102.03838v3}.

\bibitem{PS}
P.~Preu{\ss}ler and F.~L.~Schwenninger,
\emph{Implications of structured continuous maximal regularity},
arXiv:2605.12121v1, 2026.
\url{https://arxiv.org/abs/2605.12121v1}.

\bibitem{Weiss}
G.~Weiss, \emph{Admissibility of Unbounded Control Operators},
SIAM Journal on Control and Optimization \textbf{27} (1989), 527--545.
\href{https://doi.org/10.1137/0327028}{doi:10.1137/0327028}.
\end{thebibliography}

\paragraph{AI-use disclosure.}
OpenAI Codex was used to develop the argument, check mathematical and
bibliographic details, and prepare this text. An initial draft was examined
by separate AI reader and referee reviewers before the final revision.
These checks do not constitute human peer review or formal proof verification.
\end{document}
